\begin{table}[H]
    \centering \footnotesize 
    \caption{\label{tab:beliefelicitation} \textsc{Belief elicitation}}
\begin{tabularx}{\linewidth}{ X X X }
 {\bf Belief over}  & {\bf  Question}  & {\bf Possible answers}   \\ \hline
   Expected score on the PSU entrance exam, $\overline{PSU}_i^b$. &  Suppose that you will take the PSU entrance exam this year. What do you think your PSU score will be? & \begin{itemize}
    \setlength\itemsep{-0.5em}
       \item 700-850 (excellent)
       \item 600-700 (very good)
       \item 450-600 (good)
       \item 350-450 (modest)
       \item 250-350 (unsatisfactory)
       \item 150-250 (very unsatisfactory)
     %  \item I don’t know
   \end{itemize}   \\  \hline
    Expected own high school GPA, $\overline{GPA}_i^{(9-12),b}$. & Thinking of yourself, what do you think your grade point average (GPA) will be at the end of high school? (Introduce a number between 1.0 and 7.0)$^{a}$ & Free format \\ \hline
     Expected top 15\% cutoff, i.e., $85^{th}$ percentile of the GPA distribution in the school, $\bar{c}_i^{15b}$. & 
     Suppose that, in your school, there are 40 students in 12th grade. Think of the student with the 6th highest grade point average (GPA) among the 40 students. His/her GPA is in the top 15\%. What do you think is the GPA that he/she has?$^{a,b}$ & Free format \\ \hline
        %[This set of questions further elicits beliefs about the 12th student (top 30\%) and the 30th student (bottom 25\%)] & Free format \\ \hline
        Likelihood of graduating from selective college conditional on enrolling in one, $pgrad_i^b$. & If I enroll in a university (not a Technical Training Center or Professional Institute) thanks to a high PSU score, I will complete my studies.$^{c}$ & \begin{itemize} \setlength\itemsep{-0.5em}
            \item Completely certain that I will not
            \item More likely that I will not
            \item Equally likely that I will and will not
            \item More likely that I will
            \item Completely certain that I will
        \end{itemize}\\ \hline 
        
       Returns to effort in GPA and PSU productions, $\beta^{Pb}_{1i}, \beta^{Pb}_{2i}, e^{Pb}_{kink,i}, \beta^{Gb}_{1i}$. & How many hours per week do you think you need to study, between August and December 2017, to obtain a GPA/PSU score of at least $X$? [$X\in\{350, 450, 600\}$ for PSU, $X\in\{5.5,$ answer to question on top 15\% cutoff$\}$ for GPA].$^{d}$  & Free format \\ 
        \hline
   
\end{tabularx}
\begin{tablenotes}\singlespacing
\item
%\setlength{\baselineskip}{2\baselineskip}
\scriptsize
\vspace{-10pt}
\noindent \textsc{ Note.--} English translation of selected survey questions.\par
\noindent $a$. We used the Chilean term for GPA, {\itshape Notas de Ense\~nanza Media} (NEM), that refers to the grade point average across all four years of high school. Focus groups with students confirmed that this term is widely understood.\par 
\noindent $b$. Focus groups with students showed that starting by asking about the student with the highest GPA, and then asking about the student with the 6th highest, improved question comprehension. Therefore, this is how we implemented the question.\par 
\noindent $c$. Focus groups with students indicated that adding the wording `thanks to a high PSU score' was necessary to ensure students understood the question was about selective colleges, which require obtaining a PSU score above an admission cutoff, and not non-selective colleges or vocational institutions. We are confident students interpreted this question as: ``If I enroll in a selective college, I will graduate."\par 
\noindent $d$. Although the GPA question referred to the study hours required in the final semester to attain a given average high school GPA, we do not assume that respondents first calculated the additional final-semester GPA needed to attain each four-year GPA target—a cognitively demanding calculation requiring students to account for their accumulated grades and the weight of the remaining semester. We instead treat the stated GPA values as achievement targets and use the difference in reported hours to measure perceived mapping from study effort to GPA achievement.

\end{tablenotes}
\end{table}

%highest grade point average (GPA) among the 40 students. (GPA is a number between 1.0 and 7.0). What do you think is the GPA that he/she has? \par 
      %  Now think of the student with the